% !TeX root = ../papers/nc-p3.tex
\section{階数二曲線からの二変数商の検出と四辺での整合性}
\label{sec:binary-window-transfer}
$R\subset V\otimes W$を純テンソルを含まない六次元関係空間とし、
$V,W$は四次元とする。直線$D\subset V$で第一矢を割った像を
$H_D\subset(V/D)\otimes W$とする。
\begin{proposition}\label{prop:binary-quotient-rank-two-curve}
この向きで$D$を核とする二変数形式の商があるための必要十分条件は、
$\mathbb P(H_D)$内の階数一行列の軌跡に、両射影がそれぞれ
$\mathbb P(V/D),\mathbb P(W)$を線型に張る整曲線が存在することである。
このとき階数一軌跡はスキームとして有理正規五次曲線そのものであり、
二射影は滑らかな二次曲線およびねじれ三次曲線への同型である。
その逆像は$\mathbb P(R)$の階数二行列の有理正規五次曲線で、
全ての列像が$D$を含む。この曲線から$D$と二つの射影モデルが一意に戻る。
\end{proposition}
\begin{proof}
射影の核は$D\otimes W$に入り、その非零元は純テンソルなので、
$R\to H_D$は同型である。二変数形式の乗法の双対像は
$(h_{i+j})_{0\le i\le2,0\le j\le3}$の六平面である。
階数一の小行列は有理正規五次曲線のイデアルを生成する。実際、二次部分は
隣接する指数を一つずらす十本の標準二項式の空間と一致する。
こうしたずらしで同次数・同指数和の単項式を互いに移せるので、標準二項式が
パラメータ表示の核全体を生成する。
\path{computations/verify_binary_window_transfer.py}で両二次部分も比較する。

逆に指定した整曲線の関数体で、各射影の非零な座標で割り、
二組の座標関数の空間を$A,B$とする。次元は3と4で、$AB$は線型包から六次元以下である。
一般に曲線の関数体の有限次元部分空間は、滑らかな点での付値が相異なる基底を取れる。
二つの付値集合の和を取ると$\dim AB\ge\dim A+\dim B-1$で、ここでは等号である。
一般の点で消える$A$の二平面を$a\langle1,t\rangle$と書くと、$t$は非定数で、
その$B$との積は一点で消えるから五次元以下である。よって$\dim(B+tB)=5$。
$d$次元の$U$が$\dim(U+tU)=d+1$を満たせば、
$U_j=\{u:t^i u\in U,\ 0\le i\le j\}$に対し$U_{j+1}$は
$u\mapsto t^ju\bmod U_1$の核なので、各段で失う次元は高々1である。
$0\ne b\in U_{d-1}$から$U=b\langle1,t,\ldots,t^{d-1}\rangle$を得る。
従って$B$はこの形である。また$A+tA+t^2A+t^3A$は六次元で、
非定数$t$の乗法が保つ非零有限次元空間は存在しないので、各段で丁度一次元ずつ増える。
よって$A$も同じ$t$の三項の形である。曲線は二つのべき表示のグラフとなり、
その包は六平面全体である。これで必要十分性が従う。
持ち上げた行列の階数は高々2で純テンソルはないので丁度2である。
列像の射影は非退化な二次曲線を動くから、全ての列像の共通部分は丁度$D$となる。
具体的には階数二行列の曲線上で一つの共通直線を見つけ、その商の両射影の
線型包を調べれば、同じ判定を行える。
\end{proof}

\begin{corollary}\label{cor:four-window-binary-transfer}
反転を仮定しない四周期の窓で、交互の二位置で矢空間を一次元の核で割る。
各辺の向きを三次元側から見るとき前命題の条件が成立し、各頂点で左右の辺の
与える射影曲線が同一とする。このとき四窓に二変数形式の経路代数への
添字付き全射があり、長さ17には少なくとも1152次元の始点がある。
この整合性は、固定した四周期の核と埋込み曲線を用いた全射の存在と同値である。
\end{corollary}
\begin{proof}
各辺の階数一曲線の二射影は同型なので、隣接する頂点曲線間の同型を一つ与える。
一つの頂点で$\mathbb P^1$の座標を選び、この同型をたどって全ての整数位置へ
座標を運ぶ。二次または三次の完全線型系によって各座標変更は矢の線型同型に
スカラーを除き持ち上がる。スカラーの選択は関係空間を変えない。
全ての辺の関係が二変数の乗法の双対像に写るので添字付き商を得る。
逆は各辺のHankel六平面の一意な階数一曲線を読めばよい。
一周して得る$\operatorname{PGL}_2$の元は一般に恒等でなくてよい。
添字を一周ごとに区別して座標を選べるためである。四周期の座標選択まで要求する場合に限り、
射影的な一周の変換が恒等であることが必要十分となる。
\end{proof}
